\chapter{API Engineering for Crypto Venues}\label{ch:nw:api-engineering-for-crypto-venues}

A large crypto venue lets one WebSocket connection carry at most 1\,024 streams, accepts at most 300 new connections per address every five
minutes, disconnects any connection after 24 hours, and expects an answer to its ping within a minute. A firm that follows 400 symbols with
depth and trades on each has 800 streams to place, a book to rebuild whenever a message is lost, and a single budget of request weight for
the snapshots that rebuild it and the orders that trade. When one connection loses a message and the firm resynchronises every symbol on it
from full-depth snapshots, the budget runs out after 24 of them, and the last of the 80 symbols is stale for three minutes.

Chapters~16 to~19 were about where the venue is and how packets reach it. This chapter is about the interface itself: how streams are spread
over connections and addresses, how gaps are detected and repaired, how one rate budget is shared, how an \omterm{def:nw:api-engineering-for-crypto-venues:path}{order-entry path} is chosen, and
which timestamps to trust. It uses Book~3's order-book builder and rate-limit governor and Book~13's WebSocket client, and the venues'
documentation, dated.

\section{Websocket sharding}

\begin{definition}[Connection sharding]\label{def:nw:api-engineering-for-crypto-venues:shard}
\emph{Connection sharding}\index{connection sharding} is the assignment of a subscription set's streams to several connections, and of the
connections to several source addresses, so that no connection or address exceeds a venue's limits and a failure or a slow consumer affects only
part of the set.
\end{definition}

\begin{dated}{2026-09}{WebSocket limits, from venue documentation}\label{dat:nw:api-engineering-for-crypto-venues:limits}
\begin{itemize}
\item Binance spot streams: at most 1\,024 streams per connection; 300 connections per attempt every 5 minutes per IP; 5 incoming messages per
  second per connection, beyond which the connection is closed (and repeatedly disconnected IPs may be banned); connections valid 24 hours; a ping
  every 20 seconds, answered by a pong within a minute; depth events carry an event time \texttt{E}, trades a trade time \texttt{T}. A full-depth
  snapshot (up to 5\,000 levels) costs request weight 250; 1\,000 levels cost 50; the weight limit is 6\,000 a minute.
\item OKX: 3 connection requests per second per IP; 480 subscribe, unsubscribe or login requests per connection per hour; a connection without data
  for 30 seconds is closed.
\end{itemize}
\end{dated}

The venue's cap is a maximum, not a target. One connection could carry all 800 streams, but it would carry them in order: a burst on one symbol
delays every other symbol's messages behind it (chapter~1's \omterm{def:nw:networking-for-trading:hol}{head-of-line blocking} at the application level), one disconnection drops every book,
and the 24-hour reset takes the whole feed down at once. The chapter's firm caps a connection at 160 streams (80 symbols), which needs five
connections; reconnecting all five at once is far within 300 per address per five minutes, so one address suffices.

\begin{proposition}[How many connections and addresses]\label{prop:nw:api-engineering-for-crypto-venues:shards}
For $S$ streams, a cap of $c$ streams per connection (the venue's or lower) and a venue limit of $L$ new connections per address per window $W$,
a design that must be able to reconnect everything within $w \le W$ needs $\lceil S/c \rceil$ connections and $\lceil \lceil S/c \rceil / (L\,w/W)
\rceil$ addresses, if the limit is spread evenly over the window.
\end{proposition}

\begin{proof}
Each connection carries at most $c$ streams. Reconnecting $n$ connections in $w$ uses $n$ connection attempts; an address allows $L w/W$ of them in
$w$ on the even-spreading reading of the limit.
\end{proof}

\omcode{code/firm/cryptofeed/firm_cryptofeed.py}{46}{54}{The shard planner: connections at the venue's or the firm's cap, and addresses for a
reconnection within a given window.}\label{lst:nw:api-engineering-for-crypto-venues:shards}

\begin{omfigure}
\begin{tikzpicture}[font=\footnotesize,b/.style={draw,rounded corners=2pt,minimum height=0.5cm,align=center,inner sep=2pt},>=Stealth]
\node[b,fill=omDef!15,minimum height=3.4cm,minimum width=2.2cm] (s) at (0,0) {400 symbols\\ $\times$ 2 streams\\ (depth, trades)};
\foreach \k/\y in {1/1.4,2/0.7,3/0,4/-0.7,5/-1.4}{\node[b,fill=gray!10,minimum width=2.8cm] (c\k) at (4.2,\y) {connection \k: 160 streams};
  \draw[->] (s.east) -- (c\k.west);}
\node[b,fill=omThm!15,minimum height=3.4cm,minimum width=1.8cm] (a) at (8.2,0) {one source\\ address};
\foreach \k in {1,...,5}{\draw[->] (c\k.east) -- (a.west);}
\node[b,fill=omThm!25,minimum height=3.4cm] (v) at (11.2,0) {venue\\ stream\\ servers};
\draw[->,thick] (a.east) -- (v.west);
\node[font=\scriptsize,align=center,omMeth] at (4.2,-2.2) {a gap on one connection affects its 80 symbols only};
\end{tikzpicture}

\omcaption{The chapter's shard plan, schematically: 800 streams over five connections of 160, from one address. The venue would allow one connection
of 800; the firm's cap limits \omterm{def:nw:networking-for-trading:hol}{head-of-line blocking} and the damage of one disconnection.}\label{fig:nw:api-engineering-for-crypto-venues:shards}
\end{omfigure}

\section{Snapshots, deltas and sequence gaps at scale}

\begin{definition}[Per-symbol resynchronisation]\label{def:nw:api-engineering-for-crypto-venues:resync}
\emph{Per-symbol resynchronisation}\index{per-symbol resynchronisation} repairs a book after a sequence gap by fetching a snapshot for the symbols whose
own update sequence shows the gap, rather than for every symbol carried by the connection on which messages were lost.
\end{definition}

The venue's procedure is Book~3's: buffer the deltas, fetch a snapshot, drop the deltas older than it, apply the rest, and if a delta's first update
number is beyond the book's last plus one, discard the book and start again. Book~3's \texttt{firm.wsbook} implements it; fed a stream that loses one
event, it answers ``applied'', ``applied'', then ``gap''. What a gap costs is decided by how many books the firm throws away and how many snapshots its
weight budget can pay for.

\begin{proposition}[What a gap costs]\label{prop:nw:api-engineering-for-crypto-venues:gap}
If each snapshot costs $w$ of a budget of $B$ per window $T$ shared by nothing else, and $m$ books must be rebuilt, the last one returns after about
$\lceil m\,w/B \rceil - 1$ whole windows plus a round trip; with \omterm{def:nw:api-engineering-for-crypto-venues:resync}{per-symbol resynchronisation} $m$ is the number of symbols that lost messages, with
per-connection resynchronisation it is every symbol on the connection.
\end{proposition}

\begin{proof}
The governor admits $\lfloor B/w \rfloor$ snapshots per window; the remaining ones wait for the following windows.
\end{proof}

\omcode{code/firm/cryptofeed/firm_cryptofeed.py}{57}{70}{Recovery under the weight governor: each snapshot waits for the budget, and a book is stale
until its snapshot returns.}\label{lst:nw:api-engineering-for-crypto-venues:resync}

\begin{table}[ht]
\centering\small
\begin{tabular}{lrrrr}
\toprule
Snapshot depth & Weight & Per symbol (5 books) & Per connection (80 books) & Last book back \\
 & & (stale symbol-s) & (stale symbol-s) & per connection (s) \\
\midrule
5\,000 levels & 250 & 0.5 & 5\,769 & 180.1 \\
1\,000 levels & 50 & 0.5 & 11.2 & 0.2 \\
\bottomrule
\end{tabular}
\caption{Simulation: five messages lost on one connection of 80 symbols (they touch five symbols), repaired per symbol or per connection, with
snapshots at two depths under the 6\,000-weight-a-minute governor of Book~3 and a \qty{100}{\milli\second} round trip; staleness in
symbol-seconds. Data: \texttt{nw\_api.resync\_rows()}.}\label{tab:nw:api-engineering-for-crypto-venues:resync}
\end{table}

\begin{omfigure}
\begin{tikzpicture}
\begin{axis}[ybar,width=10cm,height=5.8cm,bar width=16pt,ymode=log,ymin=0.1,ymax=20000,log origin=infty,xtick=data,
  symbolic x coords={5000 levels,1000 levels},xticklabels={5\,000 levels,1\,000 levels},
  ylabel={stale symbol-seconds (log)},tick label style={font=\footnotesize},
  label style={font=\small},grid=major,grid style={gray!25},legend pos=north east,legend style={font=\footnotesize},
  table/col sep=comma,enlarge x limits=0.5]
\addplot[fill=omDef!55,draw=omDef] table[x=case,y=per_symbol]{figdata/networks/20-api-engineering-for-crypto-venues/staleness.csv};
\addplot[fill=omThm!45,draw=omThm] table[x=case,y=per_conn]{figdata/networks/20-api-engineering-for-crypto-venues/staleness.csv};
\legend{per symbol,per connection}
\end{axis}
\end{tikzpicture}

\omcaption{Simulation: stale symbol-seconds after five messages are lost on one connection, by recovery design and snapshot depth
(\cref{tab:nw:api-engineering-for-crypto-venues:resync}). Per-connection recovery with full-depth snapshots costs four orders of magnitude more.
Data: \texttt{fig\_api.py}.}\label{fig:nw:api-engineering-for-crypto-venues:stale}
\end{omfigure}

The table is the case for doing the harder thing. Per-symbol recovery needs sequence numbers per symbol (the venue's \texttt{U} and \texttt{u}) and a
book per symbol that can be discarded alone; in exchange five lost messages cost five snapshots. Per-connection recovery is simpler and, with full
depth, catastrophic: 80 snapshots of weight 250 need three and a half minutes of a 6\,000-weight budget, during which the firm has no weight left to
send orders. Snapshots of 1\,000 levels make even the careless design survivable, which is a second lesson: ask for the depth the strategy needs.

\section{Rate-limit governors across keys and addresses}

A venue's limits are per address, per account, per key, and several strategies share them. The firm needs one governor per limit that everyone
passes through (Book~3's \texttt{firm.ratelimit}), and a rule for sharing it. The chapter's allocator gives each strategy up to its demand in
proportion to a priority weight and spreads what the satisfied ones leave: with 6\,000 a minute, a risk process that needs 500 gets all of it, and
market making and arbitrage, asking 4\,000 and 3\,000 with weights 3 and 2, get 3\,300 and 2\,200. Snapshots for recovery must come out of the same
budget, which is why recovery design is a rate-budget question.

\begin{omfigure}
\begin{tikzpicture}
\begin{axis}[xbar,width=10.5cm,height=5.2cm,bar width=9pt,xmin=0,xmax=4500,ytick=data,enlarge y limits=0.25,
  yticklabels from table={figdata/networks/20-api-engineering-for-crypto-venues/budget.csv}{strategy},y dir=reverse,
  xlabel={request weight a minute},tick label style={font=\footnotesize,/pgf/number format/1000 sep={\,}},label style={font=\small},grid=major,grid style={gray!25},
  legend pos=south east,legend style={font=\footnotesize},table/col sep=comma]
\addplot[fill=gray!35,draw=gray] table[y expr=\coordindex,x=demand]{figdata/networks/20-api-engineering-for-crypto-venues/budget.csv};
\addplot[fill=omDef!55,draw=omDef] table[y expr=\coordindex,x=share]{figdata/networks/20-api-engineering-for-crypto-venues/budget.csv};
\legend{demand,share}
\end{axis}
\end{tikzpicture}

\omcaption{One 6\,000-weight minute shared by three strategies with priority weights 3, 2 and 1: the risk process's 500 is met in full, the others share
the rest in proportion. Data: \texttt{fig\_api.py}, \texttt{nw\_api.budget()}.}\label{fig:nw:api-engineering-for-crypto-venues:budget}
\end{omfigure}

\section{Choosing among order-entry paths}

\begin{definition}[Order-entry path]\label{def:nw:api-engineering-for-crypto-venues:path}
An \emph{order-entry path}\index{order-entry path} is one of the interfaces through which a venue accepts orders (REST requests, orders over a
WebSocket session, a FIX session, a venue-specific binary protocol), each with its own connection model, rate limits, latency and failure behaviour.
\end{definition}

The large venue of this chapter documents several (a REST interface, a WebSocket API, a FIX API, and market-data streams in a binary
encoding), and chapter~18's \omterm{def:nw:private-connectivity-to-crypto-venues:endpoint}{private endpoints} add routes to them. The choice follows the strategy: a persistent session removes the handshake
from each order; REST's statelessness helps recovery; FIX brings session-level sequence numbers and resend requests. Whatever the path, it draws on
the same governor, and a firm that sends orders on two paths must know whether the venue counts them against one limit or two.

\section{Redundant connections and which timestamps to trust}

\begin{definition}[Venue clock skew]\label{def:nw:api-engineering-for-crypto-venues:skew}
\emph{Venue clock skew}\index{venue clock skew} is the difference between the clock that stamps a venue's messages (its event and transaction times)
and the firm's own synchronised clock, which varies over time; it is estimated, never given.
\end{definition}

Every message carries the venue's event time; the firm stamps its arrival. The difference is the one-way delay plus the venue's clock error, and the
two cannot be separated from one side alone. Its minimum over a window, however, is the smallest delay plus the error: if the smallest delay is stable,
tracking the minimum tracks the venue's clock (chapter~4's minimum-delay filter, applied to someone else's clock).

\omcode{code/firm/cryptofeed/firm_cryptofeed.py}{91}{96}{The venue-clock estimator: the minimum of receive time minus event time over each window.}\label{lst:nw:api-engineering-for-crypto-venues:skew}

\begin{omfigure}
\begin{tikzpicture}
\begin{axis}[width=11cm,height=5.4cm,xmin=0,xmax=61,ymin=-4.6,ymax=-1.8,xlabel={minutes},
  ylabel={min(receive $-$ event) (ms)},tick label style={font=\footnotesize},label style={font=\small},grid=major,grid style={gray!25},
  legend pos=north east,legend style={font=\footnotesize},table/col sep=comma]
\addplot[omDef,thick,only marks,mark=*,mark size=1.2pt] table[x=minute,y=estimate]{figdata/networks/20-api-engineering-for-crypto-venues/skew.csv};
\addplot[omThm,thick,no markers] table[x=minute,y=truth]{figdata/networks/20-api-engineering-for-crypto-venues/skew.csv};
\legend{estimate (one-minute windows),truth: minimum delay $-$ venue offset}
\end{axis}
\end{tikzpicture}

\omcaption{Simulation: a venue clock running 3 to \qty{5}{\milli\second} ahead of the firm's over an hour, messages at ten a second with a minimum
one-way delay of \qty{0.8}{\milli\second} and exponential jitter; the one-minute minimum tracks the truth within \qty{15}{\micro\second}. Data:
\texttt{fig\_api.py}, \texttt{nw\_api.skew\_series()}.}\label{fig:nw:api-engineering-for-crypto-venues:skew}
\end{omfigure}

Redundant connections make the question sharper: two connections to the same stream deliver the same event twice, at two arrival times, and the
firm keeps the first (chapter~19's racing, for data). The venue's event time orders events across symbols and connections; the firm's arrival time
says when it could act. A strategy that uses the event time to measure its own latency must first subtract the estimated skew, or it will measure
the venue's clock.

\begin{method}[Running a crypto venue's interfaces]\label{met:nw:api-engineering-for-crypto-venues:run}
\begin{enumerate}
\item Record the venue's limits (streams, connections, messages, weights, lifetimes) with their dates, and re-read them on every change notice.
\item Shard below the caps, spread reconnections over time and addresses, and stagger the 24-hour resets.
\item Detect gaps per symbol and resynchronise per symbol, with snapshots no deeper than the strategy needs.
\item Put every request, snapshot and order through one governor per limit, and allocate it by priority.
\item Estimate each venue's clock skew continuously and timestamp everything on arrival with a synchronised clock.
\end{enumerate}
\end{method}

\section{Tutorial: four hundred symbols}\label{tut:nw:api-engineering-for-crypto-venues:tut}

\begin{tutorial}
\textbf{Goal.} Plan shards, compare recovery designs under the weight governor, share a budget and estimate a venue's clock.
\textbf{End state:} \cref{fig:nw:api-engineering-for-crypto-venues:stale,fig:nw:api-engineering-for-crypto-venues:budget,fig:nw:api-engineering-for-crypto-venues:skew}
and \cref{tab:nw:api-engineering-for-crypto-venues:resync}.
\begin{tutsteps}
\item \textbf{Shards.} \texttt{firm\_cryptofeed.plan\_shards} (\cref{lst:nw:api-engineering-for-crypto-venues:shards}) with the dated
  \texttt{LIMITS}; \texttt{nw\_api.plan()}.
\item \textbf{Gaps.} \texttt{nw\_api.gap\_demo()} feeds Book~3's \texttt{firm.wsbook} a stream with a lost event; \texttt{resync\_staleness}
  (\cref{lst:nw:api-engineering-for-crypto-venues:resync}) prices the recovery under Book~3's governor.
\item \textbf{Budget.} \texttt{allocate} shares one minute of weight.
\item \textbf{Clock.} \texttt{skew} (\cref{lst:nw:api-engineering-for-crypto-venues:skew}) on \texttt{skew\_series()}.
\end{tutsteps}
\textbf{What to change next.} Feed the recovery model Book~13's \texttt{firm.wsclient} decoded fixtures; give the governor the order traffic as
well and watch the recovery starve it.
\end{tutorial}

\section{Build: the crypto feed manager}\label{bld:nw:api-engineering-for-crypto-venues:cryptofeed}

\begin{build}
\bfield{Purpose} The interface layer to crypto venues: shards, recovery, shared budgets and venue clocks; the input of the crypto rows of chapter~29's
plan.
\bfield{Interface} \texttt{firm\_cryptofeed}: \texttt{Limits}, \texttt{LIMITS}, \texttt{plan\_shards}, \texttt{resync\_staleness}, \texttt{allocate},
\texttt{skew}; Book~3's \texttt{firm.ratelimit} and \texttt{firm.wsbook} wrapped, not edited.
\bfield{Rules} Limits are dated rows with sources; every request goes through a governor; recovery is per symbol wherever the venue's sequence numbers
allow; the venue's clock is estimated, not assumed.
\bfield{Acceptance tests} \texttt{code/firm/cryptofeed/tests/}: shard counts at the caps, recovery waiting for the next window, water-filling, and the
skew estimator by hand.
\bfield{Stretch} A live shard manager that rebalances streams by message rate; recovery by incremental snapshots where a venue offers them.
\end{build}

\begin{omsources}
\item Binance spot API documentation (WebSocket streams, local order book procedure, order-book weights); OKX API v5 documentation (WebSocket limits).
\item RFC 6455, The WebSocket Protocol; One Quant Book~3's \texttt{firm.wsbook} and \texttt{firm.ratelimit}, One Quant Book~13's \texttt{firm.wsclient}.
\end{omsources}

\section{Exercises}

\begin{exercise}[$\star$]\label{exo:nw:api-engineering-for-crypto-venues:1}
How many full-depth snapshots of weight 250 fit in a 6\,000-weight minute?
\end{exercise}

\begin{exercise}[$\star$]\label{exo:nw:api-engineering-for-crypto-venues:2}
Why cap a connection at 160 streams when the venue allows 1\,024?
\end{exercise}

\begin{exercise}[$\star$]\label{exo:nw:api-engineering-for-crypto-venues:3}
What must a client do when it receives a WebSocket ping, and what happens at this chapter's venue if it does not?
\end{exercise}

\begin{exercise}[$\star\star$]\label{exo:nw:api-engineering-for-crypto-venues:4}
Using \cref{prop:nw:api-engineering-for-crypto-venues:shards}, how many addresses does a firm need to reconnect 40 connections within 10 seconds at 300
connections per 5 minutes per address?
\end{exercise}

\begin{exercise}[$\star\star$]\label{exo:nw:api-engineering-for-crypto-venues:5}
Why does the minimum of receive time minus event time track the venue's clock, and when does it fail?
\end{exercise}

\begin{exercise}[$\star\star$]\label{exo:nw:api-engineering-for-crypto-venues:6}
A strategy measures its tick-to-order latency as its order time minus the venue's event time. What is wrong?
\end{exercise}

\begin{exercise}[$\star\star\star$]\label{exo:nw:api-engineering-for-crypto-venues:7}
\emph{Coding.} With \texttt{firm.cryptofeed}, find the largest snapshot weight at which a per-connection recovery of 80 symbols finishes within the
first minute.
\end{exercise}

\begin{exercise}[$\star\star\star$]\label{exo:nw:api-engineering-for-crypto-venues:8}
\emph{Find the flaw.} ``We subscribe to everything on one connection; the venue allows it, and fewer connections are more reliable.''
\end{exercise}

\section{Problem: Four Hundred Symbols, Five Connections}

\begin{problem}[{Weekend problem --- shards and recovery}]\label{pb:nw:api-engineering-for-crypto-venues:1}
A firm follows 400 symbols on a venue with the limits of \cref{dat:nw:api-engineering-for-crypto-venues:limits}, with a depth and a trade stream each,
at most 160 streams per connection, and one 6\,000-weight minute for snapshots and orders.

\medskip\noindent\textbf{Part I --- Shards.}
\begin{enumerate}
\item How many streams, connections and addresses does it need?
\item How many would it need at the venue's own cap?
\item How many addresses to reconnect everything within 10 seconds?
\item What does the 24-hour limit mean for the design?
\end{enumerate}

\medskip\noindent\textbf{Part II --- A gap.}
\begin{enumerate}[resume]
\item Five messages are lost on one connection. How many symbols need a snapshot per symbol, and per connection?
\item With full-depth snapshots, how long until the last book is back per connection, and how many stale symbol-seconds?
\item And with 1\,000-level snapshots?
\item What can the firm not do while a per-connection recovery runs?
\end{enumerate}

\medskip\noindent\textbf{Part III --- Budget and clocks.}
\begin{enumerate}[resume]
\item How does the allocator share 6\,000 between market making (4\,000, weight 3), arbitrage (3\,000, weight 2) and risk (500, weight 1)?
\item Where do recovery snapshots fit in that allocation?
\item How well does the one-minute minimum estimate the venue's clock in the simulation?
\item Why estimate it at all?
\end{enumerate}

\medskip\noindent\textbf{Part IV --- The verdict.}
\begin{enumerate}[resume]
\item State the \emph{named result}: the minimum connections and addresses for the subscription set under the venue's limits, and the book staleness a gap
  costs with per-symbol against per-connection resynchronisation.
\item Which design choice matters most?
\item What would change at a venue without per-symbol sequence numbers?
\item How should the shards be spread over the day's 24-hour resets?
\item What does a second connection to the same streams buy?
\item What does Book~3's governor add that a simple counter would not?
\item Where does this go in chapter~29's plan?
\item In one sentence: what does a lost message cost?
\end{enumerate}
\end{problem}

\section{Interview questions}

\begin{interviewq}[$\star$ \iqroles{developer}]\label{iq:nw:api-engineering-for-crypto-venues:1}
How do you keep a local order book correct from a snapshot and a stream of deltas?
\end{interviewq}

\begin{interviewq}[$\star\star$ \iqroles{developer}]\label{iq:nw:api-engineering-for-crypto-venues:2}
You need 2\,000 streams from a venue. How do you design the connections?
\end{interviewq}

\begin{interviewq}[$\star\star$ \iqroles{developer}]\label{iq:nw:api-engineering-for-crypto-venues:3}
Three strategies share one API key's rate limit. How do you stop one of them from starving the others?
\end{interviewq}

\begin{interviewq}[$\star\star$ \iqroles{developer, researcher}]\label{iq:nw:api-engineering-for-crypto-venues:4}
How would you estimate a venue's clock error from its market data?
\end{interviewq}

\begin{interviewq}[$\star\star$ \iqroles{developer}]\label{iq:nw:api-engineering-for-crypto-venues:5}
REST, WebSocket or FIX for order entry: how do you choose?
\end{interviewq}

\begin{interviewq}[$\star\star\star$ \iqroles{developer}]\label{iq:nw:api-engineering-for-crypto-venues:6}
Your feed handler loses messages every day at the same time. What do you look at?
\end{interviewq}
